\chapter*{초록}
\thispagestyle{plain}
본 과제는 SO-ARM101 로봇팔로 다섯 블록을 180초 이내 지정 영역으로 옮기고, 별도의 미션에서 블록을 적재한 뒤 5초간 유지하는 시스템을 목표로 하였다. 초기에는 ACT와 SmolVLA로 단일 블록 이동을 학습했으나, 시작 위치의 범위와 카메라 조건이 달라질 때 안정적인 수행을 확보하기 어려웠다. ACT와 컴퓨터 비전·역기구학(CV+IK)을 결합하는 시도에서도 파지 안정성과 시연 재수집 부담이 남아, CV+IK 기반 제어를 구현하고 실행 자료를 학습용으로 기록하도록 구성하였다.

최종 시스템은 색상·형상 기반 블록 검출, 평면 좌표 변환, IK 파지·배치 계획과 센서 기반 파지 확인을 유한상태기계로 연결하였다. 단계별 실패 분석으로 캘리브레이션, 파지 방향과 복구 동작을 보완한 뒤, 1차 미션은 30회 중 28회(93.3\%)에서 180초 내 다섯 블록 배치를 완료하였다. 재시도 포함 파지는 170회 중 148회(87.1\%) 성공했으며, 두 미션은 블록을 작업 반경 밖으로 밀어낸 뒤 시간 초과로 실패하였다. 최종 아홉 점 캘리브레이션의 적합 RMS는 5.15\,mm, 최대 leave-one-out(LOO) 오차는 13.53\,mm였다. 이는 영상 좌표와 기구학 모델 기준점 사이의 정합 오차이며, 실제 파지 정확도는 미션 수행 결과와 함께 평가하였다.

2차 미션은 5단 적재를 목표로 20회 시도하였다. 수행 도중 최대 높이 기준으로 3단에는 20회, 4단에는 17회 도달했으며 5단에는 도달하지 못했다. 재시도 포함 파지는 113회 중 94회(83.2\%) 성공하였다. 높이별 도달 횟수는 종료 시 높이나 해제 후 5초 유지 성공률을 뜻하지 않는다. 실행 자료를 자동 수집하는 Task3을 구현하여 첫 라운드에서 5개 에피소드, 2{,}111프레임을 186초에 저장하였다. 정책 재학습과 성능 비교는 수행하지 않았다.

\medskip
\noindent\textbf{주요어:} SO-ARM101, 모방학습, 컴퓨터 비전, 역기구학, 유한상태기계, 파지 검증
